\chapter{Clinical Text}
\label{chap:bt-clinical}

\begin{figure}[t]
\centering
\begin{tikzpicture}
\node[draw=ThesisNeutral, fill=ThesisPaper, rounded corners=3pt, line width=0.5pt,
      text width=0.84\textwidth, inner sep=11pt, align=left,
      font=\footnotesize, text=ThesisInk] {
{\bfseries CENTRE HOSPITALIER DE SAINT-AUBIN}\hfill Dossier n\textsuperscript{o} 4471902\\[1pt]
Service de Cardiologie --- Unité de Soins Intensifs\hfill Séjour : 03/11--08/11/2024\\[5pt]
{\color{ThesisNeutral}\rule{\linewidth}{0.3pt}}\\[5pt]
\textbf{Patient :} CARON Julien, né le 14/03/1962 (62 ans).\quad\textbf{Destinataire :} D\textsuperscript{r} S.~Lemaire.\\[6pt]
\textbf{Motif.} Douleur thoracique constrictive rétrosternale, irradiation bras G, depuis 2h.\\[5pt]
\textbf{Antécédents.} HTA. DT2 sous metformine. Tabac 30 PA. Dyslipidémie.\\[5pt]
\textbf{Examen.} TA 148/92, FC 98, SpO\textsubscript{2} 96\% AA, EVA 7/10. Pas de signe d'IVG.\\[5pt]
\textbf{Paraclinique.} ECG : sus-décalage ST V1--V4. Tropo T 4520 ng/L (N $<$14).
Coro : occlusion IVA proximale, angioplastie + stent actif, TIMI 3. ETT : FEVG 40\%.\\[5pt]
\textbf{Au total.} SCA ST+ antérieur (I21.0).\\[5pt]
\textbf{Ttt sortie.} Kardégic 75, ticagrélor 90x2, atorvastatine 80, bisoprolol 2,5, ramipril 5.\\[5pt]
\textbf{CAT.} Réadaptation cardiaque. Consultation cardio à 1 mois.\hfill D\textsuperscript{r} C.~Després
};
\end{tikzpicture}
\caption{A discharge summary for an acute myocardial infarction. The patient,
hospital, and clinicians are fictitious; the telegraphic style reproduces a real
note used only as a template. The same case returns later as scientific and lay
text.}
\label{fig:bt-clinical-note}
\end{figure}

% === P1 — HOOK ===
\Cref{fig:bt-clinical-note} is a discharge summary, written by one cardiologist
for another. It is also, to a reader of ordinary French, barely a text at all.
Articles and verbs have dropped out, findings arrive as fragments, drugs are
named by brand and dose, and the diagnosis is a code. Most of it would look, to
an outsider, like careless writing: abbreviations without explanation, units
fused to numbers, sentences that never finish. The natural reaction is to call
it bad French and to clean it up. We argue the opposite. The note is not
degraded general French but a different system, with its own grammar, written
for readers who already know almost everything it leaves out.

% === P2 ===
This intuition has a precise name. In the 1960s, Zellig Harris observed that the
language of a specialised field is not the general language with a smaller
vocabulary, but a system in its own right \citep{harris_mathematical_1968}. He
called it a \emph{sublanguage}: the sentences a community produces about a
restricted domain, governed by co-occurrence constraints tighter than, and
sometimes absent from, the language as a whole. The constraints are the point.
In a sublanguage only certain word classes may combine, so its grammar lies
partly outside ordinary usage rather than inside it as a subset. Harris and
colleagues later made the claim concrete, deriving the grammar of an immunology
subfield from its articles alone and showing that the field's information was, in
effect, its sublanguage \citep{harris_immunology_1989}.

% === P3 ===
The clinical note is the best-documented case of this idea. Building a computer
grammar of clinical English, Naomi Sager's Linguistic String Project found that
the parser needed rules for elliptical fragments that standard English does not
license, alongside word classes specific to the domain such as patient, finding,
and treatment \citep{sager_nlip_1981, sager_mlp_1987}. A contemporary panel set
out the conditions under which such a system arises: a restricted domain, a
restricted purpose, a restricted mode of communication whose time and space
limits favour compressed forms, and a community of experts who enforce the
patterns \citep{kittredge_sublanguages_1982}. A sublanguage grammar, in this
view, intersects the standard grammar without being contained in it. Sager's
further step, mapping each sentence into a fixed template of slots, is the direct
ancestor of clinical information extraction.

% === P4 ===
If the clinic runs its own grammar, so does the laboratory, and the two are not
the same. Friedman and colleagues compared clinical reports with
molecular-biology articles and described them as two distinct sublanguages, each
with its own word classes and information structure, both derivable from Harris's
theory \citep{friedman_two_2002}. The consequence matters here. ``Biomedical
text'' is not a single object: its abundant, public part, the scientific
literature, belongs to a different sublanguage from the rare, private target, the
patient note. Friedman's clinical system, \textsc{MedLEE}, shows the practical
end of the same idea, parsing radiology reports into a controlled, structured
form \citep{friedman_medlee_1994}. Abundance is therefore not the same as
relevance: more scientific text does not, on its own, supply the clinical
language we ultimately need.

% === P5 ===
The claim is measurable, not merely descriptive. \citet{temnikova_closure_2013}
quantify a sublanguage through its closure: as a corpus grows, the count of new
word types levels off toward a ceiling for a true sublanguage, whereas general
language keeps growing. Scientific corpora reach this ceiling far sooner than a
general reference corpus does. The same instruments let us test, later in this
thesis, whether clinical and scientific French are distinct sublanguages rather
than assume it.

% === P6 ===

% === P7 ===

% === P8 ===

% === P9 ===

% === P10 ===

% === P11 ===

% === P12 ===

% === P13 ===

% === P14 ===

% === P15 ===
